\documentclass[main.tex]{subfiles}

% Document
\begin{document}

\begin{abstract}
  L'esperimento internazionale DUNE si propone di far luce sui misteri ancora
  aperti della fisica dei neutrini, con l'obiettivo di misurare la
  violazione di CP, determinare l'ordinamento delle masse ed espandere le
  frontiere dell'attuale Modello Standard.\\
  Questa tesi si inserisce nello sviluppo dell'infrastruttura software di
  simulazione e ricostruzione per SAND: un rivelatore facente parte del
  complesso del \textit{Near Detector} e incaricato di monitorare continuamente
  il flusso di neutrini \textit{on-axis}.\\
  Nello specifico, il lavoro presenta lo sviluppo e l'implementazione di
  fake-reco, un modulo integrato all'interno del nuovo microFrameWork (ufw):
  un framework \textit{object-oriented} sviluppato ad hoc nell'ultimo anno
  dalla sezione INFN di Bologna.
  Il compito di questo processo è estrarre la verità Monte Carlo generata
  dagli applicativi GENIE ed EDep-Sim e trascriverla nei file di analisi
  finale \textit{Common Analysis File} (CAF), popolando in modo coerente sia i
  campi a essa dedicati sia quelli destinati alla ricostruzione.

  La scrittura di questo strumento ha permesso di comprendere a fondo la
  struttura dei CAF, la cui specifica non è stata ancora definito del tutto,
  spianando la strada all'adattamento nel ufw degli algoritmi di
  ricostruzione veri e propri.
\end{abstract}

\end{document}